% =====================================================================
%  Chapter 6: Beyond OLS — The Identification Toolkit
%  Source: BeyondOLS.tex (stub) + P3 transcription (IV, heterosk., GLS)
% =====================================================================

\chapter[Beyond OLS]{Beyond OLS: The Identification Toolkit}
\label{chap_BeyondOLS}

\section{Endogeneity: Sources and Consequences}
\label{sec_Endogeneity}

OLS is inconsistent when the key assumption $\E[X'\mathbf{e}] = 0$ fails.
This \textbf{endogeneity}\index{endogeneity} can arise from three distinct
sources.

\subsection{Omitted Variable Bias Revisited}

As shown in Section~\ref{subsec_OVB}, failing to include a variable $Z$ that
is correlated with both $X$ and $Y$ biases $\hat\beta$. The direction of bias
is determined by the sign of $(X'X)^{-1}X'Z\gamma_0$---the product of the
regression of the omitted variable on the included regressors and the effect
of the omitted variable on $Y$.

\subsection{Measurement Error}

True regressors $X^*$ are not observed; only $X = X^* + u$ is available,
where $u$ is measurement error independent of $X^*$ and $\mathbf{y}$.
Substituting:
\begin{equation}
  \mathbf{y} = X^*\beta + \mathbf{e} = (X-u)\beta + \mathbf{e}
             = X\beta + (\mathbf{e} - u\beta) = X\beta + \mathbf{v}
\end{equation}

The bias in OLS arises from $\E[X'\mathbf{v}]$:
\begin{align}
  \E[X'\mathbf{v}]
  &= \E\!\bigl[(X^*+u)'(\mathbf{e}-u\beta)\bigr] \nonumber \\
  &= \E[X^{*\prime}\mathbf{e}] + \E[u'\mathbf{e}]
     - \E[X^{*\prime}u]\beta - \E[u'u]\beta \nonumber \\
  &= -\E[u'u]\,\beta
\end{align}

This is non-zero whenever $\beta \neq 0$ and $\Var(u) > 0$, so OLS is
inconsistent. The \textbf{attenuation bias}\index{attenuation bias} pulls
$\hat\beta$ toward zero:
\begin{equation}
  \hat\beta \xrightarrow{p} \beta^* = \beta - \E\!\bigl[(X'X)^{-1}u'u\bigr]\beta
\end{equation}

\subsection{Simultaneous Equations}

When $X$ and $Y$ are jointly determined---as in a supply-and-demand system
or in a spatial model where $y_i$ and $y_j$ are mutually dependent---OLS
conflates the structural parameters of multiple equations. For example:
\begin{align*}
  q_i &= -\beta_1 p_i + e_{1i} \quad\text{(demand)} \\
  q_i &= \beta_2 p_i + e_{2i} \quad\text{(supply)}
\end{align*}
OLS on either equation alone recovers a mixture of both structural
coefficients. In each case $\hat\beta$ is \textbf{biased and
inconsistent}:\; $\partial y/\partial x = \beta + \partial e/\partial x
\neq \beta$.


\section{Instrumental Variables}
\label{sec_IV}

\subsection{The IV Estimator}

An instrument $Z$ is valid if: (a) \textbf{relevance}---$Z$ is correlated
with $X$; (b) \textbf{exclusion restriction}---$Z$ affects $Y$ only through
$X$ (not directly).

\begin{tcolorbox}[keyresult]
\begin{equation}
  \hat\beta_{IV} = (Z'X)^{-1}Z'\mathbf{y}
  \label{eq_IV}
\end{equation}
\end{tcolorbox}

Classic example \citep{Angrist1991}: month of birth and distance from
college instrument for years of schooling in a wage equation, with unobserved
ability creating endogeneity of the schooling variable.

\subsection{Two-Stage Least Squares (2SLS)}

With endogenous regressors $X_2$ and instruments $Z_2$, where $X = [X_1\;X_2]$
and $Z = [X_1\;Z_2]$ ($X_1$ instruments for itself):

\begin{enumerate}
  \item \textbf{First stage:} Regress $X$ on $Z$:\quad
    $\hat\Gamma = (Z'Z)^{-1}Z'X$,\quad $\hat X = Z\hat\Gamma$
  \item \textbf{Second stage:} Regress $\mathbf{y}$ on $\hat X$:\quad
    $\hat\beta = (\hat X'\hat X)^{-1}\hat X'\mathbf{y}$
\end{enumerate}

Scrutinising $\hat X = [\hat X_1,\hat X_2] = [P_ZX_1, P_ZX_2]
= [X_1, P_ZX_2]$: \textbf{only endogenous variables are replaced with their
fitted values}, not exogenous ones. The 2SLS estimator is:
\begin{equation}
  \hat\beta_{2SLS} = (X'P_ZX)^{-1}X'P_Z\mathbf{y}
  = \bigl[X'Z(Z'Z)^{-1}Z'X\bigr]^{-1}X'Z(Z'Z)^{-1}Z'\mathbf{y}
\end{equation}

In the just-identified case (one instrument per endogenous variable),
$\hat\beta_{2SLS} = \hat\beta_{IV}$.

\begin{tcolorbox}[keyresult]
The IV/2SLS estimator is biased in finite samples but \textbf{consistent}
under instrument validity---in contrast to OLS which is inconsistent under
endogeneity.
\end{tcolorbox}

\subsection{Spatial Lags as Instruments}

In spatial models, the endogenous variable is the spatial lag $Wy$.
A common IV strategy uses higher-order spatial lags $W^2y$, $WX$, or $W^2X$
as instruments---variables that are correlated with $Wy$ through the spatial
propagation mechanism but excluded from the structural equation. This strategy
is formalised in the SAR-IV estimator (Chapter~\ref{chap_Spatial}).


\section{Differences-in-Differences}
\label{sec_DiD}

When panel data are available and a treatment is applied at a specific point
in time, \textbf{differences-in-differences}\index{differences-in-differences}
(DiD) removes time-invariant confounders:
\begin{equation}
  Y_{it} = \alpha_i + \lambda_t + \beta(D_i \times \text{Post}_t) + \varepsilon_{it}
\end{equation}
where $\alpha_i$ are unit fixed effects, $\lambda_t$ are time fixed effects,
and $D_i \times \text{Post}_t$ is the interaction of treatment group and
post-treatment period. $\beta$ is identified under the
\textbf{parallel trends assumption}: absent treatment, treated and control
units would have followed the same time trend.

\noindent\textbf{Geographic DiD.} Administrative boundaries (census tract
borders, school district lines, city limits) often create sharp discontinuities
in treatment assignment---properties just inside versus just outside a
rezoning boundary, for example. These boundaries can serve as regression
discontinuity (RD) designs that exploit spatial proximity to identify causal
effects. See \citet{Black1999} for a classic application to school quality
capitalisation.

% TODO: Add staggered DiD and Callaway-Sant'Anna correction


\section{Heteroskedasticity and Robust Inference}
\label{sec_Heterosk}

Under heteroskedasticity, $\E[e_i^2\mid\mathbf{x}_i] = \sigma_i^2$:
\begin{equation}
  y_i = X\beta + e_i, \quad
  e_i \sim \text{id}(0, \Omega), \quad
  \Omega \equiv \operatorname{diag}(\sigma_1^2,\ldots,\sigma_n^2)
\end{equation}

OLS is still \textbf{unbiased} because unbiasedness depends only on first
moments ($\E[\mathbf{e}]=0$). But the variance is no longer $\sigma^2(X'X)^{-1}$:
\begin{align}
  V(\hat\beta)
  &= (X'X)^{-1}X'\,\E[\mathbf{e}\mathbf{e}']\,X(X'X)^{-1} \nonumber \\
  &= (X'X)^{-1}X'\Omega X(X'X)^{-1}
  \label{eq_SandwichVar}
\end{align}

This is the \textbf{sandwich (HC) estimator}\index{sandwich estimator}. OLS
standard errors computed under the homoskedastic formula $s^2(X'X)^{-1}$ are
incorrect under heteroskedasticity---typically underestimating the true
standard errors.

\begin{tcolorbox}[keyresult]
Heteroskedasticity-consistent (HC) standard errors from equation
\eqref{eq_SandwichVar} are valid under heteroskedasticity of unknown form.
Use \texttt{vcovHC()} in R or \texttt{HC3} as default.
\end{tcolorbox}


\section{Generalised Least Squares}
\label{sec_GLS}

If $\Omega$ is known (or consistently estimated), we can eliminate
heteroskedasticity by pre-multiplying the model by $\Omega^{-1/2}$:

\begin{tcolorbox}[keyresult]
\begin{equation}
  \hat\beta_{GLS} = (X'\Omega^{-1}X)^{-1}X'\Omega^{-1}\mathbf{y}
\end{equation}
\end{tcolorbox}

GLS is the Best Linear Unbiased Estimator (BLUE) when $\Omega$ is known.
In practice $\Omega$ must be estimated (Feasible GLS), introducing estimation
error but retaining consistency.

\noindent\textbf{Spatial GLS.} In spatial models the error covariance
$\Omega$ has a specific parametric structure determined by the spatial weight
matrix $W$ and the error dependence parameter $\lambda$. The Spatial Error
Model (SEM, Chapter~\ref{chap_Spatial}) is a GLS model with
$\Omega = [(I-\lambda W)'(I-\lambda W)]^{-1}\sigma^2$.


\section{Panel Data Methods}
\label{sec_Panel}

Panel data follow the same units (individuals, firms, cities) over multiple
time periods. The key advantage is the ability to control for time-invariant
unobserved heterogeneity.

\subsection{Fixed Effects}

The \textbf{within estimator}\index{fixed effects estimator} removes
unit-specific means:
\begin{equation}
  (y_{it} - \bar{y}_i) = (x_{it} - \bar{x}_i)'\beta + (e_{it} - \bar{e}_i)
\end{equation}

This eliminates any unit-level time-invariant confounder
$\alpha_i$---whether or not it is observed. The fixed effects estimator is
consistent under the weaker assumption that $\E[x_{it}e_{it}] = 0$ (strict
exogeneity within unit), not $\E[x_{it}\alpha_i] = 0$.

\subsection{Random Effects}

The random effects estimator treats unit effects $\alpha_i$ as random
variables drawn from a distribution, rather than fixed parameters to be
estimated. It is more efficient than fixed effects but requires the stronger
assumption that $\alpha_i \ind X_{it}$---unit effects are uncorrelated with
regressors. The \textbf{Hausman test} discriminates between the two:
if fixed effects and random effects estimates differ significantly, the random
effects assumption is rejected.

\subsection{Two-Way Fixed Effects and DiD}

The two-way FE model adds time fixed effects $\lambda_t$:
\begin{equation}
  y_{it} = \alpha_i + \lambda_t + X_{it}'\beta + e_{it}
\end{equation}
Under parallel trends, this is numerically equivalent to the
differences-in-differences estimator described above. Recent work by
\citet{Callaway2021} shows that the standard two-way FE estimator can produce
misleading estimates when treatment adoption is staggered across units; they
propose a cohort-based correction.

% TODO: Add Callaway-Sant'Anna worked example with urban policy data

\renewcommand\bibname{Further reading}
\begin{subappendices}
\bibliographystyle{econometrica}
\bibliography{Literature/OverLord}
\IfFileExists{Literature/BeyondOLS.bbl}{\chapter{Beyond OLS}\label{chap_BeyondOLS}
\section{Measurement error and identification}
\section{Instrumental variables (IV), Two-Stage Least Squares (2SLS)}




\renewcommand\bibname{Further reading}
\begin{subappendices}
\bibliographystyle{econometrica}
\bibliography{Literature/OverLord}
\IfFileExists{Literature/BeyondOLS.bbl}{\chapter{Beyond OLS}\label{chap_BeyondOLS}
\section{Measurement error and identification}
\section{Instrumental variables (IV), Two-Stage Least Squares (2SLS)}




\renewcommand\bibname{Further reading}
\begin{subappendices}
\bibliographystyle{econometrica}
\bibliography{Literature/OverLord}
\IfFileExists{Literature/BeyondOLS.bbl}{\chapter{Beyond OLS}\label{chap_BeyondOLS}
\section{Measurement error and identification}
\section{Instrumental variables (IV), Two-Stage Least Squares (2SLS)}




\renewcommand\bibname{Further reading}
\begin{subappendices}
\bibliographystyle{econometrica}
\bibliography{Literature/OverLord}
\IfFileExists{Literature/BeyondOLS.bbl}{\input{Literature/BeyondOLS.bbl}}{\typeout{Need to run bibtex}}
\end{subappendices}
}{\typeout{Need to run bibtex}}
\end{subappendices}
}{\typeout{Need to run bibtex}}
\end{subappendices}
}{\typeout{Need bibtex on BeyondOLS}}
\end{subappendices}
